\chapter{Newton Polygon}

\section{Definitions}

In this section we will determine a way to check for zeroes of a power series via the construction
of Newton polygons. We consider power series over a non-archimedean valued field $R$ where we denote
the \emph{additive} valuation by $v$ and the associated multiplicative valuation by $\lvert \cdot \rvert$.
The typical example would be $\Q_p$.

These are formally given by the following definition.

\begin{definition}
    The \emph{Newton polygon} of $f = \sum a_i x^i$ is the lower boundary of the convex hull
    of the set of points $(i,\nu (a_i))$, where we ignore points with infinite $p$-adic valuation.
    Here by convex hull we mean the smallest convex polygon that contains all points.
\end{definition}

Given a Newton polygon we can give the following information.

\begin{definition}
    \begin{enumerate}
        \item The \emph{slopes} of the line segments appearing in the Newton polygon are called the
        slopes of $f(x)$.
        \item The \emph{length} of a slope is the length of the projection of the line segment with
        that slope onto the $x$-axis.
        \item The \emph{breaks} are the values of $i$ which are vertices of the Newton polygon.
    \end{enumerate}
\end{definition}

It turns out when working with Newton polygons the slopes and lengths are the information we care about;
and moreover given a starting point using this we can reconstruct all breaks. This provides motivation
for the definition we give in Lean.

\begin{definition}
  A Newton polygon is the following information:
  \begin{itemize}
    \item a natural number $n$, determining the number of slopes;
    \item a (potentially infinite) list of slopes that are strictly increasing except for the final
      slope (which may only be equal to the previous);
    \item a (potentially infinite) list of lengths; and
    \item a starting point.
  \end{itemize}
\end{definition}

It is clear that given a Newton polygon, in the formal sense, we can construct the following information;
however the reverse also holds by using the comment above the definition and the fact the slopes are
(strictly) increasing give the convexity condition.
%% Perhaps I could elaborate on this further?

We now want to construct a way to obtain this information given a power series, $f = \sum_{i=0}^\infty a_n x^i \in R[x]$.
Towards this we consider the following algorithm:

\begin{enumerate}
    \item Plot the points $\{(i,\nu (a_i) )\}$.
    \item Start with the vertical half-line down from the point with the smallest $i$.
    \item Rotate the line until one of the following happens:
    \subitem The line simultaneously hits infinitely many of the points we have plotted. In this
        case don't break the line and the polygon is complete.
    \subitem The line reaches a position where it can be rotated no further without leaving behind
        some points. In this case, the half lines forms the final segment of the polygon.
    \subitem The line hits a finite number of points. In this case, break at the last point that was
        hit and repeat the procedure.
\end{enumerate}

It should be clear that by construction this constructs the lower convex hull of the set of points
$\{(i,\nu (a_i) )\}$ and taking the slopes and lengths of the obtained segments we get our corresponding
definition in Lean.

As mentioned, the study of Newton polygons is deeply connected to studying the zeroes of $f(x)$,
thus we will assume $f(0) \neq 0$. That is, we factor out any powers of $x$ dividing $f$, which
means our initial vertical line starts on the $y$-axis.

Moreover, we will also assume $a_0 = 1$. This is because if we factor
$f(x) = \sum_{i = 0}^\infty a_i x^i$ as $f(x) = a_0 \cdot g(x)$, where
$g(x) = 1 + \sum_{i=1}^\infty \frac{a_i}{a_0}x^i$, it is not hard to see that the resulting Newton
polygon of $g(x)$ is the Newton polygon of $f(x)$ translated down by $\nu (a_0)$.
This fact follows from observing $\nu (\frac{a_i}{a_0}) = \nu(a_i) - \nu (a_0)$.

\section{Properties}

Let $K$ be a $p$-adic field, either a finite extension of $\mathbb{Q}_p$ or $\mathbb{C}_p$.
%% can this be generalised to any nonarch-valued field (probably?)

We want to consider a power series and consider its norms $\lvert f(x) \rvert_c$ for many different
$c$.

It is easy to see that
\[
c_1 > c_2 \implies \lvert f(x) \rvert_{c_1} \geq \lvert f(x) \rvert_{c_2}.
\]

Let us consider $f = 1 + a_1 x + \cdots a_n x^n$ a polynomial and consider its first segment. If this
segment has slope $m$, it connects the point $(0,0)$ to some other point $(i, mi)$ (i.e. slope $m$,
length $i$). Then we get the following properties:
\begin{itemize}
  \item $\lvert a_j \rvert \leq p^{-mj} = (p^{-m})^j$  for all $j$;
  \item $\lvert a_i \rvert = p^{-mi}$; and
  \item $\lvert a_j \rvert < p^{-mj}$ if $j > i$.
\end{itemize}

If we let $c = p^m$, we can read these conditions as
\begin{itemize}
  \item $\lvert f(x) \rvert_c = 1$; and
  \item $i$ is the largest integer such that $\lvert f(x) \rvert_c = \lvert a_i \rvert c^i.$
\end{itemize}

That is we see that $f$ is $i$-distinguished. In particular, this means, by Weierstrass preparation
that if $i$ is less than the degree of $f$ then, $f$ is divisible by a polynomial of degree $i$.

%% Note I actually need to formalise the statements about weirstrass div and prep that focuses on
%% if they are only polynomials now!

\begin{definition}
  A polynomial $f(x)$ is pure if its Newton polygon has only one slope. If this slope is $m$, we say
  $f$ is pure of slope $m$.
\end{definition}

\begin{proposition}
  Irreducible polynomials are pure.
\end{proposition}

\begin{proof}
  A break at $(i,mi)$ means a factor of degree $i$; hence, if there is a break at $i \neq 0, n$
  (so not pure) and hence the polynomial is reducible.
\end{proof}

\begin{proposition}
  A polynomial, of degree $n$, will be pure of slope $m$ if and only if it has the property that for
  $c = p^m$
  \[
  \lvert f(x) \rvert_c = \lvert b_n \rvert c^n = 1,
  \]
  i.e. it is $n$-distinguished.
\end{proposition}

\begin{theorem}
  Let $f(x) = 1 + a_1 x + \cdots + a_n x^n \in K[x],$ and assume the newton polygon of $f$ has its
  first break at $(i,mi)$. Then there exist polynomials $g(x), h(x) \in K[x]$ satisfying:
  \begin{itemize}
    \item $f(x) = g(x) h(x)$;
    \item $g(x)$ has degree $i$ and is pure of slope $m$; and
    \item $h(x)$ has no zeros in the closed ball of radius $p^m$ around 0.
  \end{itemize}
\end{theorem}

\begin{proof}
  Put everything above together.
\end{proof}

\begin{lemma}
  Let $f(x) = 1 + a_1 x + \cdots a_n x^n$, and assume the first break of the Newton polygon occurs
  at $(i,mi)$. Let $c$ be any positive real number less than $p^m$. Then we have
  \[
    \lvert f(x) \rvert_c = 1 \quad \text{and} \quad \lvert f(x) - 1 \rvert_c < 1.
  \]
\end{lemma}

\begin{proof}
  If $\lvert f(x) - 1 \rvert_c < 1$, then we must have $\lvert f (x) \rvert_c = 1$ by the usual all
  triangles are isosceles argument. Thus, we only need to prove the inequality.
  A line through the origin with slope $m_1 < m$ passes below all points on the polygon, touching it
  only at $(0,0)$. This means that $v_p(a_j) > m_1 j$ for every $j > 0$. Or in terms of absolute
  values $\lvert a_j \rvert < p^{-m_1 j}.$ Since $a_0 = 1$, the zeroth coefficient of $f(x) - 1$ is
  zero, and therefore also $\lvert a_0 - 1 \rvert < 1.$ It follows that $\lvert f(x) - 1 \rvert_c < 1$
  for $c < p^{m_1}$.
\end{proof}

\begin{lemma}
  Let $f(x) = 1 + a_1 x + \cdots + a_n x^n$, and assume that the first break of the newton polygon
  occurs at the point $(i,mi)$. Let $c$ be any positive real number less than $p^m$. Then $f$ has no
  zeros in the closed ball in $\mathbb{C}_p$ of radius $c$ around 0.
\end{lemma}

\begin{proof}
  The previous lemma says that $\lvert f(x) - 1 \rvert_c < 1$. From that, it follows for any $x$ such
  that $\lvert x \rvert \leq c$, we have $\lvert f(x) - 1 \rvert < 1$, which means $f(x) \neq 0$.
  Thus, $f(x)$ has no zeros in $\overline{B}(0,c)$.
\end{proof}

%% the reason we have C_p and just just Q_p bar is that we eventually want to talk about
%% convergence of power series... which requires completeness

%% Therefore, in Lean this can be changed. But the later will have to be in C_p
%% For generality this will be algebraic closure of an ultrametric normed ring (field?)
%% We then want its completion later on (see how C_p is defined by Maria Ines)

\begin{proposition}
  Let $f(x) = 1 + a_1 x + a_2 x^2 + \cdots a_n x^n$, and assume that the first break of the newton
  polygon occurs at the points $(i,mi)$. Then $f(x)$ has no roots with absolute value less than $p^m$
  and has exactly $i$ roots (counting multiplicity, in $\mathbb{C}_p$) with absolute value $p^m$.
\end{proposition}

\begin{proof}
  If the first break is at $(i,mi)$ then
  \begin{itemize}
    \item if $c < p^m$, $f(x)$ has no roots in the closed ball of radius $c$; and
    \item $f(x)$ factors as the product of a pure polynomial $g(x)$ of slope of $m$ and a polynomial
      which has no roots in the closed ball of radius $p^m$.
  \end{itemize}
  Any root of $g$ is a root of $f$, and therefore has no roots of absolute value less than $p^m$. On
  the other hand, if $\alpha_1, \cdots, \alpha_i$ are the roots of $g$, then we must have
  \[
    g(x) = (1 - \alpha_1^{-1}x) (1 - \alpha_2^{-1}x) \cdots (1 - \alpha_i^{-1}x),
  \]
  and hence the $i$-th coefficient is equal to $(\prod_i \alpha_i)^{-1}$. Since $g$ is pure, this
  must have valuation $mi$, and as $\lvert (a_j) \rvert \geq p^m$, it follows that all $a_j$ have
  absolute value $p^m$.
\end{proof}

\begin{theorem}
  Let $f(x) = 1 + a_1 x + \cdots + a_n x^n$ be a polynomial, let $m_1, \cdots, m_r$ be the slopes
  of its Newton polygon (in increasing order), and let $i_1, \cdots, i_r$ be their corresponding
  projected lengths. Then, for each $k$, $1 \leq k \leq r$, $f$ has exactly $i_k$ roots
  (in $\mathbb{C}_p$, counting multiplicity) of absolute value $p^{m_k}$.
\end{theorem}

\begin{proof}
  We have already seen this for the first segment. So let us move onto the second. Note this suffices
  to complete the proof as the same argument will yield the result for consecutive segments.

  Therefore, let us assume there are breaks at $(i_1,mi)$ and $(i_1 + i_2, m_1 i + m_2 i_2)$. We
  will set $k := i_1 + i_2$. Similar to the first slope we can conclude the following:
  \begin{itemize}
    \item $\lvert a_k \rvert = p^{-m_2 i_2 - m_1 i_1} = p^{-m_2 k} p^{(m_2 - m_1)i}$;
    \item for $i_1 \leq j \leq k$, $\lvert a_j \rvert \leq p^{-m_2 (j - i_1) - m_1 i_1}$; and
    \item for $j < i_1$ and $ j > k$, we have $\lvert a_j \rvert < p^{-m_2 (j - i_1)- m_1 i_1}.$
  \end{itemize}
  In other words, we get for $c = p^{m_2}$ that $\lvert f \rvert_c = p^{(m_2 - m_1)i_1}$, i.e. $f$
  is $k$-distinguished. Therefore, we can go through a completely analogous proof to before to get
  $f$ has exactly $k$ roots in the closed ball of radius $p^{m_2}$, $i_1$ of which have absolute
  value $p^{m_1}$ and $k-i_1 = i_2$ of which have absolute value $p^{m_2}$,
\end{proof}

With this we can now move onto considering power series and give the key theorem.

\begin{theorem} \label{RadofCon}
    Let $m$ be the supremum of the slopes appearing in the Newton polygon of a series
    $f(x) = 1 + a_1 x + a_2 x^2 + \cdots$. Then the radius of convergence of $f$ is $p^m.$
\end{theorem}

\begin{proposition}\label{NP}
    Let $f(x) = 1 + a_1 x + a_2 x^2 + \cdots$ be a power series. Let $m_1, m_2, \dots, m_k$ be the
    first $k$ slopes of the Newton polygon of $f(x)$, and assume that $f(x)$ converges on the
    closed ball of radius $c = p^{m_k}$ (from previous section, this means it is a restricted power
    series for parameter $c$ if $K$ is complete). Let $N$ be the $x$-coordinate of the right
    endpoint of the $k$-th segment of the Newton polygon. Then, there exists a polynomial
    $g(x)$ of degree $N$ and a power series $h(x)$ such that
    \begin{enumerate}
        \item $f(x) = g(x) h(x)$,
        \item $\| f(x) - g(x) \|_c < 1$,
        \item $h(x)$ converges on the closed ball of radius $c$,
        \item $\| h(x) - 1\|_c <1$, and
        \item the Newton polygon of $g(x)$ is the same as the portion of the Newton polygon of
        $f(x)$ contained in the region $0 \leq x \leq N$.
    \end{enumerate}
\end{proposition}

This then allows us to give our connection to zeroes of the power series.

\begin{corollary}
    Let $f(x) = 1 + a_1 x + a_2 x^2 + \cdots$ be a power series which converges on the closed ball
    of radius $c = p^m$. Let $m_1,\dots, m_k$ be the slopes of the Newton polygon of $f(x)$ which
    are less than or equal to $m$, and let $i_1,\dots,i_k$ be their lengths. Then, for each $j$, the
    power series $f(x)$ has $i_j$ zeros with absolute value $p^{m_j}$, and there are no other zeros
    in the closed ball of radius $p^m.$
\end{corollary}

\begin{proof}[Proof of proposition~\ref{NP}]
    We prove this by induction on $k$:\\
    If $k=1$, write the break point as $(N, m_1 N)$. We know all points are on or above the line
    $y = m x$, and the ones where $j > N$ are strictly above it. This translates to
    \begin{itemize}
        \item $| a_j | (p^j)^{m_1} \leq 1$ for all $j$,
        \item $| a_N | (p^N)^{m_1} = 1$, and
        \item $| a_j | (p^{m_1})^j$ < 1 if $j > N.$
    \end{itemize}
    This says that $\|f(x) \|_{p^{m_1}} = 1$ and that the maximum is last realised at the degree
    $N$. Thus, we can apply the Weierstrass preparation theorem to conclude that there is a
    polynomial of degree $N$ and a power series $h(x)$ satisfying
    \[
    \| h(x) - 1 \|_{p^m_1} < 1 \quad \|f(x) - g(x)\|_{p^m_1} < \|f(x) \|_{p^m_1} \quad \text{and} \quad f(x) = g(x) h(x).
    \]
    Combining the middle equation with $\| f(x) \|_{p^{m_1}} < 1$ means
    $\| f(x) - g(x) \|_{p^{m_1}} < 1$.
    Now, using facts about Newton polygons of polynomials, it follows that the zeros of $f(x)$ in
    the closed ball of radius $p^m$ coincide with the those of $g(x)$. That is, we must have the
    Newton polygon of $g(x)$ coincides with the Newton polygon of $f(x)$ on $0 \leq x \leq N$.
    Finally, we have $f(x)$ converges on the closed ball of radius $p^{m_1}$, by Theorem
    \ref{RadofCon}, and thus so does $h(x).$
    \\
    By induction, assume we have $f(x) = g_{k-1}(x)h_{k-1}(x)$ for a polynomial $g_{k-1}(x)$ and
    power series $h_{k-1}(x)$, such that the Newton polygon of $g_{k-1}(x)$ coincides with the
    first $k-1$ segments of the polygon for $f(x)$ and $h_{k-1}(x)$ has no zeroes in the closed ball
    of radius $p^{m_{k-1}}$.
    \\
    Let us now consider the $k$-th segment. By definition, the $k$-th segment ending at $x = N$
    means that for any $i > N$, $(i, \nu_p (a_i))$ lies above the line of slope $m_k$ through the
    point $(N, \nu_p (a_N))$. This means that, writing $c = p^{m_k}$, $\| f(x) \|_{p^{c}} =
    \lvert a_N \rvert c^N$ and $\lvert a_n \rvert c^N > \lvert a_i \rvert c^i $ for any $i > N$.
    Thus, we can apply the $p$-adic Weierstrass preparation theorem to get a polynomial $g(x)$.
    It is then easy to see $g_1 (x) \mid g(x)$, and that its Newton polygon coincides with the
    relevant portion of the Newton polygon of $f(x).$
\end{proof}

\begin{proof}[Proof of corollary]
    Assume we know this statement for polynomials, and the proposition says that the relevant
    part of the Newton polygon of $f(x)$ \emph{is} the Newton polygon of the polynomial $g(x)$.
    Since $g(x)$ is a factor of $f(x)$ and the quotient $h(x)$ clearly has no zeroes in the closed
    ball of radius $p^m$, the conclusion follows.
\end{proof}
